\DocumentMetadata{
  pdfversion=2.0,
  pdfstandard=ua-2,
  lang=en-US,
  tagging=on,
  tagging-setup={math/setup=mathml-SE,tabsorder=structure}
}
\documentclass[11pt,letterpaper]{article}
\usepackage[margin=0.68in,includefoot]{geometry}
\usepackage{newcomputermodern}
\usepackage{amsmath,amscd,mathtools}
\providecommand{\square}{\mdlgwhtsquare}
\usepackage[hidelinks]{hyperref}
\hypersetup{
  pdftitle={Analysis First-Year Exam, May 2019, Part 1},
  pdfauthor={University of Florida Department of Mathematics},
  pdfsubject={Graduate mathematics examination},
  pdfdisplaydoctitle=true,
  bookmarksopen=true
}
\setcounter{secnumdepth}{0}
\setlength{\parindent}{0pt}
\setlength{\parskip}{0.28em}
\setlength{\emergencystretch}{3em}
\setlength{\leftmargini}{2.15em}
\setlength{\leftmarginii}{2.65em}
\setlength{\leftmarginiii}{3em}
\renewcommand{\labelenumi}{\textbf{\theenumi.}}
\pagestyle{empty}
\raggedbottom
\allowdisplaybreaks
\newenvironment{examproblems}[1]{%
  \begin{enumerate}%
  \setcounter{enumi}{#1}%
  \setlength{\topsep}{0.35em}%
  \setlength{\partopsep}{0pt}%
  \setlength{\itemsep}{0.48em}%
  \setlength{\parsep}{0.18em}%
}{\end{enumerate}}
\newenvironment{examsubparts}{%
  \begin{enumerate}%
  \setlength{\topsep}{0.18em}%
  \setlength{\partopsep}{0pt}%
  \setlength{\itemsep}{0.16em}%
  \setlength{\parsep}{0.1em}%
}{\end{enumerate}}
\newenvironment{examinstructions}{%
  \par\begingroup\itshape
}{%
  \par\endgroup\vspace{0.18em}
}
\makeatletter
\renewcommand\section{\@startsection{section}{1}{0pt}%
  {-1.25ex plus -.25ex minus -.1ex}%
  {0.55ex plus .1ex}%
  {\normalfont\large\bfseries}}
\renewcommand{\@maketitle}{%
  \newpage\null\vskip -1.2em%
  {\footnotesize\bfseries
    \href{https://www.ufl.edu/}{University of Florida}, \href{https://math.ufl.edu/}{Department of Mathematics}\par}%
  \vskip 0.18em%
  \tagstructbegin{tag=Title}%
    {\LARGE\bfseries\href{https://gma.math.ufl.edu/exams/analysis-first-year/}{Analysis First-Year Exam}, May 2019, Part 1\par}%
  \tagstructend%
  \vskip 0.35em%
  \tagmcbegin{artifact}%
    \hrule height 1pt%
  \tagmcend%
  \vskip 0.55em%
}
\makeatother
\title{Analysis First-Year Exam, May 2019, Part 1}
\author{}
\date{}
\begin{document}
\maketitle
\thispagestyle{empty}
\begin{examinstructions}
Do exactly 4 problems. Work must be presented in a neat and logical fashion in order to receive credit. Do not leave any gaps. When a theorem is used in a proof it must be precisely stated.
\end{examinstructions}
\begin{examproblems}{0}
\item
Let \((a_n)\) and \((b_n)\) be bounded real sequences, with \(\lim a_n=a\). Prove that 
\[
\limsup(a_n+b_n)=a+\limsup b_n.
\]
\item
Let \((X,d)\) be a metric space and \(E\subseteq X\) a nonempty subset. Consider the function \(d_E:X\to[0,+\infty)\) defined by 
\[
d_E(x)=\inf_{y\in E}d(x,y).
\]
\begin{examsubparts}
\item[\textnormal{i)}]
Prove \(d_E\) is uniformly continuous on~\(X\).
\item[\textnormal{ii)}]
Describe, with proof, the set of points at which \(d_E=0\).
\end{examsubparts}
\item
Let~\(X\) be a compact metric space. Let \(\{F_\alpha\}_{\alpha\in A}\) be a (nonempty) collection of nonempty closed subsets of~\(X\), and suppose that \(\{F_\alpha\}\) is totally ordered by inclusion (this means that for every \(\alpha,\beta\in A\), either \(F_\alpha\subseteq F_\beta\) or \(F_\beta\subseteq F_\alpha\)). Prove that the intersection \(\bigcap_{\alpha\in A}F_\alpha\) is nonempty.
\item
Let \(\{U_\alpha\}_{\alpha\in A}\) be a family of connected subsets of a metric space~\(X\). Suppose that \(U_\alpha\cap U_\beta\ne\varnothing\) for each pair \(\alpha,\beta\in A\). Prove that \(\bigcup_{\alpha\in A}U_\alpha\) is connected.
\item
Let \(f:\mathbb{R}\to\mathbb{R}\) be a differentiable function, and suppose that \(f'\) is monotone. Must \(f'\) be continuous? Prove, or give a counterexample.
\end{examproblems}
\end{document}
